\chapter{LLM-Based Oracle Category Normalization}
\label{app:oracle-category-normalization}

The natural-language oracle returned 46 distinct category labels for BBQ and 236 for
MedQA. Orthographic normalization alone reduced the MedQA vocabulary only to 222 labels
and did not reduce the BBQ vocabulary. To obtain interpretable co-grouping heatmaps, an
\textbf{OpenAI Codex large language model} assigned every raw label to exactly one
dataset-specific high-level category. This is an analytical post-processing step and
not an output of the FELICS auxiliary oracle.

The assignment followed a fixed rubric. First, capitalization, Unicode dash variants,
spacing, grammatical number, and obvious synonymous forms were treated as equivalent.
Second, each label was assigned according to its primary semantic role. Composite or
ambiguous labels were assigned to the most specific applicable family; no label was
duplicated across families. Third, the resulting mapping was frozen and applied
identically to both evaluated models and all examples. The large language model saw
only the raw category label during assignment: concept texts, group memberships,
model outputs, and calculated co-grouping rates were not used. Consequently, the
normalization cannot be interpreted as an independently validated clinical or social
science taxonomy. The complete auditable mapping is stored in
\path{data/oracle_category_mapping.csv}, and the augmented concept-level data are
stored in \path{data/oracle_assignments_normalized.csv}.

Let $\pi(x,y)$ denote an unordered pair. For normalized categories $a$ and $b$,
the displayed co-grouping rate is
\begin{equation}
  r_{ab}=
  \frac{\sum_{q}\sum_{i<j}
    \mathbb{1}[\pi(A_{qi},A_{qj})=\pi(a,b)]\,
    \mathbb{1}[G_{qi}=G_{qj}]}
  {\sum_{q}\sum_{i<j}
    \mathbb{1}[\pi(A_{qi},A_{qj})=\pi(a,b)]},
\end{equation}
where $q$ indexes a question/model run, $A_{qi}$ is the normalized category of concept
$i$, and $G_{qi}$ is its oracle group. Pairs are never formed across questions or runs. A cell is undefined when
its denominator is zero and is shown in gray.

\section*{BBQ category families}

\begingroup\small
\begin{longtable}{@{}>{\raggedright\arraybackslash}p{0.31\textwidth}p{0.62\textwidth}@{}}
\hline
\textbf{High-level category} & \textbf{Assignment scope} \\
\hline
\endfirsthead
\hline
\textbf{High-level category} & \textbf{Assignment scope} \\
\hline
\endhead
Demographics & Age, gender, sexual orientation or sexuality, appearance, and general demographic descriptors. \\
Relationships \& family & Relationships and their status or type, family structure or status, and intimacy. \\
Actions \& behavior & Actions, activities, behavior, interaction, and action-related statements or objects. \\
Topics \& communication & Topics of conversation or interest, communication, questions, and information sources. \\
Intentions \& attitudes & Intentions, preferences, attitudes or opinions, interests, concerns, and uncertainty. \\
Emotions \& reactions & Emotions, feelings, reactions, and mixed intention/emotion labels. \\
Setting \& time & Locations, environments, time, and mixed location/event labels. \\
Events \& situations & Events, situations, experiences, and challenges. \\
Identity \& roles & Names and social or narrative roles. \\
Other details & Non-specific detail labels that could not be assigned more narrowly. \\
\hline
\end{longtable}
\endgroup

\section*{MedQA category families}

\begingroup\small
\begin{longtable}{@{}>{\raggedright\arraybackslash}p{0.31\textwidth}p{0.62\textwidth}@{}}
\hline
\textbf{High-level category} & \textbf{Assignment scope} \\
\hline
\endfirsthead
\hline
\textbf{High-level category} & \textbf{Assignment scope} \\
\hline
\endhead
Demographics & Age, sex or gender, race/ethnicity, and socioeconomic status. \\
Symptoms & Presenting, systemic, organ-specific, pain, rash, sleep, and other symptom descriptions. \\
Timing \& course & Onset, duration, chronology, progression, follow-up, and temporal relationships. \\
Medical history & General medical, allergy, congenital, collateral, and disease-control history. \\
Family, social \& lifestyle & Family and social history, behavior, diet, occupation, living situation, lifestyle, and substance use. \\
Reproductive history & Obstetric, gynecologic, menstrual, pregnancy, prenatal, gestational, and reproductive information. \\
Medication \& adherence & Medication exposure or history, medication changes, adherence, and compliance. \\
Treatment \& response & Management, interventions, treatment plans or timelines, and clinical response. \\
Vitals \& anthropometrics & Vital signs, anthropometric measurements, body weight, obesity, and derived measurements. \\
Physical examination & General and system-specific examinations, appearance, maneuvers, measurements, and examination findings. \\
Laboratory \& pathology & Laboratory panels, microbiology, urinalysis, biopsy, pathology, and specimen findings. \\
Imaging & Imaging and radiographic findings, including infarct localization or type. \\
Diagnostic tests & Non-imaging diagnostic studies such as ECG, audiometry, polysomnography, and ophthalmologic testing. \\
Anatomy \& localization & Anatomical structure, location, laterality, distribution, and localized mass descriptors. \\
Neurologic function & Neurologic examination, sensory or motor deficits, nerve involvement, and related function. \\
Mental health \& cognition & Mental status, cognition, psychiatric or psychological information, affect, and suicide risk. \\
Diagnosis \& reasoning & Diagnoses, provisional or differential diagnoses, and diagnostic considerations. \\
Care context & Reason for visit, care setting or utilization, prior care, prevention goals, and patient concerns. \\
Risk, exposure \& trauma & Risk factors, environmental exposure, trauma, travel, and potentially relevant recent activity. \\
\hline
\end{longtable}
\endgroup
